\section{From lexical knowledge to observed production}
\label{sec:ch3:chapter-introduction}

The first empirical study in this thesis asks how a model's vocabulary growth can be compared with children's when the comparison is based on what each learner produces. Chapter~\ref{chap:litreview} established that such a comparison requires a specified relation between input, learning, and measured behavior. Expressive vocabulary makes this requirement concrete. A learner may encounter a word many times without producing it in a short observation, and a model may generate a plausible word form without establishing that it understands its meaning. The present chapter develops an explicit measurement framework for this problem and uses it to examine the lexical trajectories of predictive language learners.

Vocabulary is informative because it supports comparisons at several levels simultaneously: how many target words appear in production, how their availability changes with experience, and whether frequent and infrequent words follow different trajectories. The MacArthur--Bates Communicative Development Inventories (CDI) and the developmental data assembled in Wordbank provide a substantial basis for studying early expressive vocabulary through caregiver reports \citep{fenson2007macarthur,frank2017wordbank}. Naturalistic child transcripts provide a complementary record of observed productions. Together, these sources make lexical development a useful setting in which to ask whether learning from bounded quantities of linguistic input yields patterns resembling those found in children.

The measures nevertheless capture different aspects of vocabulary. A caregiver inventory summarizes whether a child is reported to say particular words across everyday situations. A transcript records the forms observed during sampled interactions. Model generation supplies another sample of forms, whose composition depends on the trained distribution and the decoding procedure. Throughout this chapter, \emph{lexical knowledge} refers to the broader capacity to represent and use words, \emph{expressive vocabulary} to the words available for production, and \emph{observed production} to the output actually recorded under a specified procedure. The benchmark operationalizes the last of these as evidence about the second; it does not independently measure the full lexical knowledge of either learner.

\meCDI{} makes this link explicit through three coordinated choices. It scales training corpus sizes to estimated quantities of childhood input, sets generation budgets using estimates of monthly child production, and constructs content-word evaluation sets aligned in lexical category and frequency. A count threshold calibrated on child productions links corpus observations to the CDI reference trajectory. These choices address a central ambiguity in vocabulary comparison: a larger vocabulary count can result from more learning, more opportunities to produce words, or a test list containing easier items. The benchmark brings these factors into the experimental design rather than leaving them implicit.

The study compares character-level LSTMs and Transformers trained on child-directed speech or audiobook transcripts. It asks whether increasing the available training data reproduces the child production trajectory, whether the answer converges across vocabulary coverage, lexical diversity, and non-word rate, and whether decoding temperature or lexical frequency accounts for the observed differences. Age labels index estimated input and output budgets; the experiment does not simulate biological maturation or one child's continuous learning history.

The clearest result is that the evaluated models do not jointly reproduce the child reference patterns. Target-word coverage increases with input size, but its trajectory differs from the accelerated growth in the child corpus; low-frequency words are particularly sensitive to model and sampling conditions. Lexical diversity and non-word rates also depend on architecture and corpus, making them complementary constraints rather than interchangeable measures of progress. The following study presents the benchmark and these comparisons in detail. The chapter discussion then considers what they establish about data efficiency and how they motivate the subsequent test of retrieval on syntactic contrasts.
